\documentclass[11pt]{article}

\usepackage[a4paper,margin=29mm]{geometry}
\usepackage[T1]{fontenc}
\usepackage[utf8]{inputenc}
\usepackage{lmodern}
\usepackage{microtype}
\usepackage{amsmath,amssymb,amsthm,mathtools}
\usepackage{booktabs}
\usepackage{array}
\usepackage{enumitem}
\usepackage[hidelinks]{hyperref}

\newtheorem{theorem}{Theorem}[section]
\newtheorem{proposition}[theorem]{Proposition}
\newtheorem{corollary}[theorem]{Corollary}
\newtheorem{lemma}[theorem]{Lemma}
\theoremstyle{definition}
\newtheorem{definition}[theorem]{Definition}
\newtheorem{problem}[theorem]{Problem}
\theoremstyle{remark}
\newtheorem{remark}[theorem]{Remark}

\newcommand{\PG}{\mathrm{PG}}
\newcommand{\F}{\mathbb{F}}
\newcommand{\cC}{\mathcal{C}}
\newcommand{\cH}{\mathcal{H}}
\newcommand{\cU}{\mathcal{U}}
\newcommand{\cX}{\mathcal{X}}
\newcommand{\IC}{I_{\cC}}
\newcommand{\IH}{I_{\cH}}
\newcommand{\rhoC}{\rho_{\cC}}
\newcommand{\tmin}{t_2(2,q)}

\title{Arcs complete outside a prescribed conic\\
\large A prescribed-hole defect identity and a certified classification over $\F_{16}$}
\author{Tavis Rudd}
\date{July 2026}

\begin{document}
\maketitle

\begin{abstract}
Let $\cC$ be a nonsingular conic in $\PG(2,q)$.  We study arcs
$A\subseteq\PG(2,q)\setminus\cC$ whose secants cover every point outside
$A\cup\cC$, and write $\rhoC(q)$ for their minimum size.  Starting from the
classical first two equations for the secant-index distribution of an arc,
we prove an exact defect identity for an arbitrary prescribed set of points
that may remain uncovered, together with equality and stability criteria.
Its conic specialization gives
\[
 q^2-k\le
 \binom{k}{2}(q-1)-\frac{6}{\lfloor k/2\rfloor}\binom{k}{4}
 -\frac{\IC(A)}{\lfloor k/2\rfloor}
\]
for every $\cC$-complete $k$-arc, together with an exact equality criterion
and a quantitative stability estimate.  In particular, for every prime power $q$,
\[
 \rhoC(q)\ge \sqrt{2q}+\frac32-\frac8{\sqrt{2q}}.
\]
Conversely, if $H$ is any prescribed set and an ordinary complete $b$-arc
satisfies $b|H|<q^2+q+1$, projective averaging moves it off $H$; the conic
case gives $\rhoC(q)\le \tmin$ whenever $\tmin\le q$.  We also give tangent
and parity constraints in even characteristic.  A sharp finite-field
evaluation dichotomy packages the quadratic obstruction for arbitrary
feature maps, and the arc--MDS dictionary turns the defect identity into an
exact syndrome-leader collision identity.  Directly checkable witnesses
give
\[
 \rhoC(8)=\rhoC(9)=\rhoC(11)=6,
 \qquad \rhoC(16)=9.
\]
A classification of all $2633$ projective eight-arcs over $\F_{16}$ proves
the stronger statement that no nonzero quadratic can contain an eight-arc's
ordinary-uncovered locus while avoiding the arc.  A supplementary verifier
checks the witnesses, while a companion Lean formalization proves the theorem
chain and kernel-checks the local classification, quadratic obstruction,
projective reduction, and exact $q=16$ conclusion.  For the six-point
$q=11$ witness it certifies the no-conic premise which, by the classical
normal-rational-curve dictionary, makes the associated $[6,3,4]_{11}$ MDS
code projectively non-GRS; it also proves covering radius three, the exact
coset and extension spectra, and the icosahedral simultaneous-extension complex.
\end{abstract}

\section{Introduction}

An arc in a projective plane is a set of points no three of which are
collinear.  An arc is complete when it cannot be enlarged while remaining
an arc; equivalently, every point outside it lies on a secant determined by
two of its points.  Complete arcs and the closely related problem of finding
small saturating sets have a substantial literature; see, for example,
Ball~\cite{Ball1997}, Kim--Vu~\cite{KimVu2003}, and Nagy~\cite{Nagy2018}.

Fix a nonsingular conic $\cC\subset\PG(2,q)$.  We impose the arc condition
but allow the points of $\cC$ to remain uncovered.

\begin{definition}\label{def:relative}
An arc $A\subseteq\PG(2,q)\setminus\cC$ is \emph{$\cC$-complete} if every
point of
\[
 \PG(2,q)\setminus(A\cup\cC)
\]
lies on a secant of $A$.  Define
\[
 \rhoC(q)=\min\{|A|:A\text{ is a $\cC$-complete arc}\}.
\]
For an arbitrary arc disjoint from $\cC$, its uncovered locus is
\[
 \cU_{\cC}(A)=
 \PG(2,q)\setminus
 \left(A\cup\cC\cup\bigcup_{\{a,b\}\in\binom A2}ab\right).
\]
Thus $A$ is $\cC$-complete exactly when $\cU_{\cC}(A)=\varnothing$.
\end{definition}

The minimum in Definition~\ref{def:relative} is attained.  Indeed, in the
finite family of arcs disjoint from $\cC$, any member of maximum cardinality
is maximal under insertion and hence is $\cC$-complete.  The same argument
works for every prescribed hole set.

Equivalently, $A$ is maximal among arcs in the point set
$\PG(2,q)\setminus\cC$.  All nonsingular conics in $\PG(2,q)$ are
projectively equivalent, so the numerical value of $\rhoC(q)$ does not
depend on the chosen conic.

The problem differs from two nearby standard notions.  A saturating set
need not be an arc and has no prescribed exceptional locus.  An
almost-complete subset of a conic, in the sense of
Bartoli--Davydov--Marcugini--Pambianco~\cite{BartoliEtAl2018}, consists of
selected points \emph{on} the conic whose chords cover the complementary
points.  Ng--Wild's arcs covering a line~\cite{NgWild2001} prescribe a
one-dimensional set that must be covered, rather than an exceptional set
that may remain uncovered.  Here the selected arc is required to avoid the
conic.

There is also classical work on \emph{complete exterior sets} of a conic:
sets of $(q+1)/2$ exterior points whose pairwise joins are exterior lines; see
Blokhuis--Seress--Wilbrink~\cite[pp.~143, 146]{BlokhuisSeressWilbrink1992},
who report that, up to isomorphism, one of the two $q=11$ configurations is a
six-arc.  This is the configuration described thirty-five years earlier by
Edge~\cite[\S\S29--32]{Edge1956}, who named it a \emph{Clebsch hexagon} after
its real-plane antecedent.  This is a
different condition from prescribed-hole completeness: it constrains how
the selected points meet the conic, whereas Definition~\ref{def:relative}
requires their secants to cover the complementary off-conic locus.  The
$q=11$ witness below has all of its secants exterior to the conic, making
this literature directly relevant even though the paper's completeness
condition is different.  In uncovered-locus notation, the classical
exterior-set condition gives only $\cC(\F_{11})\subseteq U(A)$.  The reverse
inclusion used below follows by combining Dye's ten-Brianchon-point result
with the elementary chord-defect count; it is not stated in the BSW paper.

The point index with respect to an arc, and the first two double-counting
equations for its distribution, are classical; compare
Hirschfeld~\cite[Chapter~9]{Hirschfeld1998} and Ball~\cite{Ball1997}.  The
use of secant counts to bound the smallest complete $(k,n)$-arcs is also
classical; for a recent explicit lower bound see
Alabdullah--Hirschfeld~\cite{AlabdullahHirschfeld2019}.  The
main observation of this note is that, after splitting those equations over
a prescribed exceptional set and its complement, the inequalities normally
used to control overlap have an exact remainder.  This produces:

\begin{enumerate}[label=(\roman*),leftmargin=2em]
\item an exact prescribed-hole defect identity and equality criterion;
\item a quantitative stability bound for near equality;
\item a strengthened lower bound with additive asymptotic term $3/2$;
\item a projective-averaging upper-bound transfer from ordinary complete
      arcs;
\item characteristic-two nucleus constraints;
\item an exact evaluation-avoidance dichotomy and a coding-theoretic form of
      the prescribed-hole defect; and
\item small finite-field results whose upper bounds are independently
      verifiable from explicit coordinates.
\end{enumerate}

Thus the classical scaffolding is the secant-index equations, the
Lunelli--Sce $\sqrt{2q}$ scale, the arc--MDS dictionary, and standard
saturating-set terminology.  The new claims isolated here are the
prescribed-hole parameter, the exact remainder after the exceptional-set
split (including its equality and stability consequences), and the use of
the ordinary-uncovered locus as an evaluation obstruction.  The defect
identity is deliberately an exact accounting theorem derived from the two
classical moments; its content is the identified nonnegative remainder and
the consequences obtained by retaining it, not a claim that the underlying
moments are new.

A targeted literature search across complete and saturating arcs,
almost-complete conic subsets, prescribed-line coverage, arcs and quadrics,
and the published classification of arcs in $\PG(2,16)$ did not locate
Definition~\ref{def:relative} or the exact defect identity below in this
form.  This observation is not a priority certificate: the proofs are
elementary consequences of classical index equations, and terminology
involving relative saturation or holes may occur elsewhere.

\section{The classical secant equations}
\label{sec:classical}

Throughout the remainder of the paper, $q\ge3$ and $k\ge3$.  Let $\Pi$ be a projective plane of order $q$, let $A$ be a $k$-arc in
$\Pi$, and let $r_A(x)$ be the number of secants of $A$ through
$x\in\Pi\setminus A$.  We usually write $r(x)$.  Put
\[
 N=\binom{k}{2},
 \qquad
 m=\left\lfloor\frac{k}{2}\right\rfloor.
\]

\begin{lemma}\label{lem:max-index}
For every $x\notin A$, one has $r(x)\le m$.
\end{lemma}

\begin{proof}
Distinct secants through $x$ cannot share an endpoint in $A$.  Indeed, two
lines through the same pair $x,a$ coincide.  Hence the secants through $x$
use pairwise disjoint pairs of points of $A$, of which there are at most
$\lfloor k/2\rfloor$.
\end{proof}

\begin{proposition}[Classical first and second index equations]
\label{prop:moments}
For every $k$-arc $A$ in a projective plane of order $q$,
\begin{align}
 \sum_{x\notin A}r(x)
   &=\binom{k}{2}(q-1), \label{eq:first-moment}\\
 \sum_{x\notin A}\binom{r(x)}2
   &=3\binom{k}{4}. \label{eq:second-moment}
\end{align}
\end{proposition}

\begin{proof}
There are $\binom{k}{2}$ secants.  Each contains $q+1$ points, exactly two
of them in $A$, and therefore contributes $q-1$ incidences with points
outside $A$.  This proves~\eqref{eq:first-moment}.

For~\eqref{eq:second-moment}, count unordered pairs of distinct secants
through a common external point.  By Lemma~\ref{lem:max-index}, their four
endpoints in $A$ are distinct.  Conversely, a four-subset of $A$ has three
partitions into two unordered pairs.  Each partition gives two secants,
which meet in a unique point because the plane is projective.  Their
intersection is outside $A$: otherwise one of the two secants would contain
three points of the arc.  Thus every four-subset contributes exactly three.
\end{proof}

Only the projective-plane axioms and the common line size $q+1$ enter these
proofs.  In particular, the second equation need not hold in a general
finite linear space where two lines may be disjoint.

\section{An exact defect identity for prescribed holes}
\label{sec:defect}

The next theorem is stated for an arbitrary exceptional point set.  The
conic enters only in later specializations.

Let $\cH\subseteq\Pi\setminus A$ be any set of points allowed to remain
uncovered, and define
\begin{align*}
 V_{\cH}(A)&=\Pi\setminus(A\cup\cH),\\
 \cX_{\cH}(A)&=\{x\in V_{\cH}(A):r(x)>0\},\\
 \cU_{\cH}(A)&=V_{\cH}(A)\setminus\cX_{\cH}(A),\\
 \IH(A)&=\sum_{y\in\cH}r(y).
\end{align*}
Thus $\cX_{\cH}(A)$ is the covered required locus, and $\cU_{\cH}(A)$ is
the uncovered required locus.

Splitting Proposition~\ref{prop:moments} over $V_{\cH}(A)$ and $\cH$ gives
\begin{align}
 \sum_{x\in V_{\cH}(A)}r(x)
   &=N(q-1)-\IH(A), \label{eq:split-first}\\
 \sum_{x\in V_{\cH}(A)}\binom{r(x)}2
 +\sum_{y\in\cH}\binom{r(y)}2
   &=3\binom{k}{4}. \label{eq:split-second}
\end{align}

\begin{theorem}[Prescribed-hole defect identity]
\label{thm:defect}
Let $A$ be a $k$-arc in a projective plane of order $q$, let
$\cH\subseteq\Pi\setminus A$, and put $m=\lfloor k/2\rfloor$.  Define
\[
 \Delta_{\cH}(A)=
 N(q-1)-\frac{6}{m}\binom{k}{4}
 -\frac{\IH(A)}{m}-|\cX_{\cH}(A)|.
\]
Then
\begin{equation}
\boxed{
 m\Delta_{\cH}(A)=
 \sum_{x\in\cX_{\cH}(A)}(r(x)-1)(m-r(x))
 +\sum_{y\in\cH}r(y)(m-r(y)).
}
\label{eq:defect}
\end{equation}
In particular, $\Delta_{\cH}(A)\ge0$.
\end{theorem}

\begin{proof}
Only points of $\cX_{\cH}(A)$ contribute to the sum in
\eqref{eq:split-first}.  Multiplying the definition of $\Delta_{\cH}(A)$
by $m$ and using~\eqref{eq:split-first} gives
\begin{align*}
 m\Delta_{\cH}(A)
 &=m\sum_{x\in\cX_{\cH}(A)}r(x)
   +(m-1)\IH(A)-m|\cX_{\cH}(A)|
   -6\binom{k}{4}.
\end{align*}
By~\eqref{eq:split-second},
\[
 6\binom{k}{4}=
 2\sum_{x\in\cX_{\cH}(A)}\binom{r(x)}2
 +2\sum_{y\in\cH}\binom{r(y)}2.
\]
Substitution and termwise simplification yield
\begin{align*}
 m(r-1)-2\binom r2&=(r-1)(m-r),\\
 (m-1)r-2\binom r2&=r(m-r),
\end{align*}
which proves~\eqref{eq:defect}.  Every term is nonnegative by
Lemma~\ref{lem:max-index}.
\end{proof}

\begin{corollary}[Coverage and uncovered-locus bounds]
\label{cor:coverage}
Under the hypotheses of Theorem~\ref{thm:defect},
\begin{equation}
 |\cX_{\cH}(A)|\le
 N(q-1)-\frac{6}{m}\binom{k}{4}-\frac{\IH(A)}m,
\label{eq:coverage}
\end{equation}
and
\begin{equation}
 |\cU_{\cH}(A)|\ge
 |V_{\cH}(A)|-N(q-1)
 +\frac{6}{m}\binom{k}{4}+\frac{\IH(A)}m.
\label{eq:uncovered}
\end{equation}
Equality holds in either inequality if and only if
\[
 r(x)\in\{1,m\}\quad(x\in\cX_{\cH}(A)),
 \qquad
 r(y)\in\{0,m\}\quad(y\in\cH).
\]
\end{corollary}

\begin{proof}
The two inequalities are respectively $\Delta_{\cH}(A)\ge0$ and its
rearrangement using
$|\cU_{\cH}(A)|=|V_{\cH}(A)|-|\cX_{\cH}(A)|$.
The equality statement follows from the vanishing of every nonnegative term
in~\eqref{eq:defect}.
\end{proof}

For later use, we record the complete-outside consequence before imposing
any geometry on the holes.

\begin{corollary}[Arbitrary prescribed holes]
\label{cor:arbitrary-holes}
Let $|\cH|=h$.  If $A$ is a $k$-arc disjoint from $\cH$ and complete outside
$\cH$, then
\begin{equation}
 q^2+q+1-k-h\le
 \binom{k}{2}(q-1)-\frac6m\binom{k}{4}-\frac{\IH(A)}m.
\label{eq:arbitrary-holes}
\end{equation}
\end{corollary}

\begin{proof}
Completeness gives
$|\cX_{\cH}(A)|=|V_{\cH}(A)|=q^2+q+1-k-h$.
Substitute this equality into~\eqref{eq:coverage}.
\end{proof}

The exact remainder also quantifies near equality.

\begin{corollary}[Quantitative stability]
\label{cor:stability}
Assume $m\ge3$ and put
\begin{align*}
 M&=\{x\in\cX_{\cH}(A):2\le r(x)\le m-1\},\\
 J&=\{y\in\cH:1\le r(y)\le m-1\}.
\end{align*}
Then
\begin{equation}
 (m-2)|M|+(m-1)|J|\le m\Delta_{\cH}(A).
\label{eq:stability}
\end{equation}
\end{corollary}

\begin{proof}
For $2\le r\le m-1$,
$(r-1)(m-r)\ge m-2$.  For $1\le r\le m-1$,
$r(m-r)\ge m-1$.  Apply these estimates to~\eqref{eq:defect}.
\end{proof}

Thus small defect forces almost every multiply covered required point to
have the maximum possible multiplicity $m$, while almost every hole point
met by a secant also has multiplicity $m$.  This is the precise stability
content behind the equality heuristic.

\section{Coding interpretation}
\label{sec:coding}

We record the standard arc--code dictionary because it makes the defect
identity more informative.  Choose nonzero column representatives of a
$k$-arc $A\subset\PG(2,q)$ and form a $3\times k$ matrix $H_A$.  For
$k\ge4$, the code
\[
 C_A=\{c\in\F_q^k:H_Ac^{\mathsf T}=0\}
\]
is a $[k,k-3,4]_q$ MDS code: every three columns are independent, the
columns span $\F_q^3$, and any four columns support a dependence.  Conversely,
a projective parity-check system for a code with these parameters is a
$k$-arc.  This correspondence and its coset-weight consequences are
classical; see Davydov--Marcugini--Pambianco~\cite{DavydovEtAl2021}.
We call such a code \emph{projectively GRS} when, after permuting and scaling
columns and applying a projective coordinate change, its projective
parity-check columns lie on a normal rational curve; in $\PG(2,q)$ this is a
nonsingular conic.  This is the standard normal-rational-curve/GRS
correspondence~\cite{StormeThas1991}.

\begin{proposition}[Syndrome form of relative completeness]
\label{prop:syndrome-dictionary}
Let $s\ne0$ have projective direction $x=[s]\in\PG(2,q)$.  Then:
\begin{enumerate}[label=(\roman*),leftmargin=2em]
 \item the minimum weight of a word $c$ with $H_Ac^{\mathsf T}=s$ is one,
       two, or three according as $x\in A$, $x\notin A$ lies on a secant of
       $A$, or $x$ is ordinarily uncovered;
 \item if the minimum weight is two, the number of weight-two leaders is
       exactly $r_A(x)$;
 \item if the minimum weight is three, the number of weight-three leaders is
       exactly $\binom{k}{3}$; and
 \item adjoining $x$ as one more projective parity-check column preserves
       the MDS property exactly when $x$ is ordinarily uncovered.
\end{enumerate}
Consequently, $A$ is complete outside a prescribed set $\cH$ exactly when
all its projective distance-three syndrome directions lie in $\cH$ (and
$A\cap\cH=\varnothing$).  These directions are deep-hole directions only
when covering radius three has separately been established.
\end{proposition}

\begin{proof}
A weight-one representation is a scalar multiple of a selected column.  A
weight-two representation uses a unique unordered column pair whose
projective span is a secant through $x$; independence of the pair makes its
two coefficients unique.  If neither case occurs, any three columns form a
basis and represent $s$, so the distance is three.  For each three-subset
the representation is unique, and none of its coefficients can vanish;
hence there are exactly $\binom{k}{3}$ leaders.  The extension statement is
the same three-column independence test.  The leader count also follows as
the $d=4$ case of the general farthest-coset formula in
Davydov--Marcugini--Pambianco~\cite[Theorem~7.7]{DavydovEtAl2021}; the
deep-hole/MDS-extension formulation is standard in the Reed--Solomon setting
as well~\cite{Kaipa2017}.
\end{proof}

Thus $r_A(x)$ is not merely an incidence index: it is the exact number of
minimum weight-two leaders of every affine syndrome on direction $x$.
All $q-1$ nonzero scalar multiples of a projective syndrome direction have
the same distance and leader count: multiplication by the scalar bijects
their coefficient words without changing Hamming support.  This is the
normalization behind every projective-to-affine factor of $q-1$ below.
The first moment counts leaders, the second counts collisions between
distinct leader supports, and Theorem~\ref{thm:defect} is an exact
leader-collision remainder identity after prescribed directions are removed.
In particular, Corollary~\ref{cor:arbitrary-holes} is also a length
obstruction for projective codimension-three MDS codes whose distance-three
syndrome directions are confined to the prescribed projective set.  This is
a reformulation of the geometric theorem, not a claim that the arc--MDS or
deep-hole dictionaries are new.

\section{Specialization to a conic}
\label{sec:conic}

Let $\cC$ be a nonsingular conic in $\PG(2,q)$, and let
$A\subseteq\PG(2,q)\setminus\cC$ be a $k$-arc.  Since $|\cC|=q+1$, there
are $q^2-k$ points outside $A\cup\cC$.  Write
\[
 \IC(A)=\sum_{y\in\cC}r(y)
       =\sum_{\ell\text{ an $A$-secant}}|\ell\cap\cC|.
\]

\begin{corollary}[Conic-loss and uncovered-locus inequalities]
\label{cor:conic-bound}
For every such arc,
\begin{equation}
 |\cU_{\cC}(A)|\ge
 q^2-k-\binom{k}{2}(q-1)
 +\frac{6}{m}\binom{k}{4}+\frac{\IC(A)}m.
\label{eq:conic-uncovered}
\end{equation}
If $A$ is $\cC$-complete, then
\begin{equation}
\boxed{
 q^2-k\le
 \binom{k}{2}(q-1)-\frac{6}{m}\binom{k}{4}
 -\frac{\IC(A)}m.
}
\label{eq:conic-complete}
\end{equation}
\end{corollary}

The first correction in~\eqref{eq:conic-complete} accounts for unavoidable
concurrence among secants; the last accounts for secant incidence spent on
the exempt conic.

Define
\begin{align*}
 L_1(q)&=\min\left\{k\ge3:q^2-k\le\binom{k}{2}(q-1)\right\},\\
 L_2(q)&=\min\left\{k\ge3:
 q^2-k\le\binom{k}{2}(q-1)
 -\frac{6}{\lfloor k/2\rfloor}\binom{k}{4}\right\}.
\end{align*}
Then
\begin{equation}
 \rhoC(q)\ge L_2(q)\ge L_1(q).
\label{eq:L2}
\end{equation}
The definition of $L_2(q)$ involves only integer arithmetic after
multiplication by $\lfloor k/2\rfloor$.

For calculation, the corrected capacity has the parity-dependent forms
\begin{equation}
 \binom{k}{2}(q-1)-\frac{6}{\lfloor k/2\rfloor}\binom{k}{4}
 =
 \begin{cases}
 \displaystyle
 \frac{k-1}{2}\bigl(k(q-1)-(k-2)(k-3)\bigr),&k\text{ even},\\[6pt]
 \displaystyle
 \frac{k}{2}\bigl((k-1)(q-1)-(k-2)(k-3)\bigr),&k\text{ odd}.
 \end{cases}
\label{eq:parity-capacity}
\end{equation}
For comparison, some values are
\[
\begin{array}{c|rrrrrrrr}
q&5&8&9&11&13&16&17&19\\ \hline
L_1(q)&4&5&5&6&6&7&7&7\\
L_2(q)&4&6&6&6&7&8&8&8
\end{array}
\]
Thus the second-moment correction alone proves, for example, that a
$\cC$-complete arc in $\PG(2,16)$ has at least eight points.
At $q=5$, the matching lower bound $L_2(5)=4$ has a four-frame witness:
the companion Clebsch analysis checks that its ordinary-uncovered locus is
the full rational point set of a nonsingular conic.  A queued follow-up will
import that witness into the present formal semantics and promote the derived
target $\rhoC(5)=4$ into the verified small-value theorem; the companion owns
the broader classification for arcs of sizes four through seven.

\begin{theorem}[Explicit additive lower bound]
\label{thm:asymptotic}
For every prime power $q$,
\begin{equation}
 \boxed{\displaystyle
 \rhoC(q)\ge\sqrt{2q}+\frac32-\frac8{\sqrt{2q}}.}
\label{eq:asymptotic-lower}
\end{equation}
Consequently,
\begin{equation}
 \liminf_{q\to\infty}
 \bigl(\rhoC(q)-\sqrt{2q}\bigr)\ge\frac32.
\label{eq:liminf}
\end{equation}
\end{theorem}

\begin{proof}
Let $k$ be the size of a $\cC$-complete arc and put $s=\sqrt{2q}$.
The elementary first-moment inequality
$q^2-k\le\binom{k}{2}(q-1)$ implies $k\ge s$ for $q\ge2$:
otherwise $k<s\le q$, while $k^2<2q$ gives
$\binom{k}{2}<q$, so its right side is smaller than $q(q-1)$ whereas
$q^2-k\ge q(q-1)$.
If $k\ge s+2$, then~\eqref{eq:asymptotic-lower} is immediate.  Otherwise
write $k=s+a$, so $0\le a<2$ and $s\ge2$.

Since $m=\lfloor k/2\rfloor\le k/2$,
\[
 \frac6m\binom{k}{4}
 \ge\frac{12}{k}\binom{k}{4}
 =\frac{(k-1)(k-2)(k-3)}2.
\]
Dropping the nonnegative conic term in~\eqref{eq:conic-complete} therefore
gives the necessary inequality
\[
 q^2-k\le
 \frac{k-1}{2}\bigl(k(q-1)-(k-2)(k-3)\bigr).
\]
After substituting $q=s^2/2$ and $k=s+a$, the right side minus the left
side is exactly
\[
 \frac{2a-3}{4}s^3
 +\frac{a^2-7a+10}{4}s^2
 +\frac{-3a^2+10a-8}{2}s
 +\frac{-a^3+5a^2-8a+6}{2}.
\]
For $0\le a\le2$, the last three coefficients are at most
$5/2$, $1$, and $3$, respectively.  Since $s\ge2$, their total
contribution is at most $4s^2$.  The displayed difference is nonnegative,
and hence
\[
 0\le \frac{2a-3}{4}s^3+4s^2.
\]
Dividing by $s^2$ gives $a\ge3/2-8/s$, which is
\eqref{eq:asymptotic-lower}.  Letting $q$ tend to infinity gives
\eqref{eq:liminf}.
\end{proof}

\begin{remark}[The scale of the conic term]
\label{rem:scale}
Every secant meets $\cC$ in at most two points, so
\[
 0\le\IC(A)\le2\binom{k}{2}.
\]
At the relevant scale $k\asymp\sqrt q$, the term $\IC(A)/m$ is therefore
$O(\sqrt q)$, while the universal overlap correction
$6\binom{k}{4}/m$ is $\Theta(q^{3/2})$.  Consequently, improving the lower
bound on $\IC(A)$ within~\eqref{eq:conic-complete} can be useful for finely
balanced finite cases, but cannot improve the additive constant $3/2$ in
Theorem~\ref{thm:asymptotic}.  Any stronger asymptotic result needs an
additional mechanism, not merely a spectral estimate for $\IC(A)$.
\end{remark}

\section{An upper-bound transfer from complete arcs}
\label{sec:transfer}

Let $t_2(2,q)$ denote the smallest size of an ordinary complete arc in
$\PG(2,q)$.

\begin{theorem}[Projective averaging]
\label{thm:transfer}
Let $H\subseteq\PG(2,q)$ and let $B$ be an ordinary complete $b$-arc.  If
\begin{equation}
 b|H|<q^2+q+1,
\label{eq:transfer-general}
\end{equation}
then some projective image of $B$ is disjoint from $H$ and complete outside
$H$.
\end{theorem}

\begin{proof}
Choose a uniformly random projectivity $g\in\operatorname{PGL}(3,q)$.
Point transitivity gives
\[
 \mathbb E|gB\cap H|=\frac{b|H|}{q^2+q+1}<1.
\]
Since the intersection size is a nonnegative integer, it is zero for some
$g$.  Projectivities preserve ordinary completeness, so $gB$ is complete
outside $H$.
\end{proof}

\begin{corollary}[Conic transfer]
\label{cor:conic-transfer}
If $t_2(2,q)\le q$, then
\begin{equation}
 \rhoC(q)\le t_2(2,q).
\label{eq:transfer}
\end{equation}
More generally, every complete $b$-arc with $b\le q$ has a projective image
disjoint from the prescribed conic.
\end{corollary}

\begin{proof}
Take $H=\cC$, so $|H|=q+1$.  For integral $b\le q$ one has
$b(q+1)<q^2+q+1$, and Theorem~\ref{thm:transfer} applies.
\end{proof}

Kim and Vu proved that, for all sufficiently large $q$, every projective
plane of order $q$ has a complete arc of size at most
$\sqrt q(\log q)^c$ for an absolute constant $c$~\cite{KimVu2003}.  This is
less than $q$ for large $q$, so Corollary~\ref{cor:conic-transfer} gives the
following immediate consequence.

\begin{corollary}\label{cor:kim-vu}
For some absolute constant $c$ and all sufficiently large prime powers $q$,
\[
 \rhoC(q)\le\sqrt q(\log q)^c.
\]
\end{corollary}

The transfer does not exploit the conic and is unlikely to be sharp.  It
nevertheless supplies a uniform upper bound without having to preserve the
arc condition during a separate random covering construction.

\section{Even characteristic and the nucleus}
\label{sec:nucleus}

Assume $q$ is even and let $\nu$ be the nucleus of $\cC$.  The nucleus is
outside $\cC$, all tangents to $\cC$ pass through $\nu$, and every line
through $\nu$ is one of these $q+1$ tangents.

\begin{proposition}[The nucleus belongs to the arc]
\label{prop:nucleus-in}
Let $A$ be a $k$-arc disjoint from $\cC$ with $\nu\in A$.  Exactly $k-1$
$A$-secants are tangent to $\cC$.  In particular,
\[
 \IC(A)\ge k-1,
 \qquad
 \IC(A)\equiv k-1\pmod2.
\]
If $A$ is $\cC$-complete, then
\begin{equation}
 q^2-k\le
 \binom{k}{2}(q-1)-\frac6m\binom{k}{4}-\frac{k-1}{m}.
\label{eq:nucleus-in}
\end{equation}
\end{proposition}

\begin{proof}
For every $a\in A\setminus\{\nu\}$, the line $\nu a$ is tangent to
$\cC$, giving $k-1$ tangent secants.  No line joining two points of
$A\setminus\{\nu\}$ can be tangent: every tangent contains $\nu$, which
would place three points of $A$ on one line.  Hence these are exactly the
tangent $A$-secants.

Every tangent contributes one to $\IC(A)$, while every other line
contributes zero or two.  This proves the lower bound and parity statement.
Inequality~\eqref{eq:nucleus-in} follows from
Corollary~\ref{cor:conic-bound}.
\end{proof}

\begin{proposition}[The nucleus lies outside the arc]
\label{prop:nucleus-out}
Let $A$ be a $\cC$-complete arc with $\nu\notin A$.  Then
\[
 1\le r(\nu)\le m,
\]
the tangent $A$-secants are exactly the secants through $\nu$, and
\[
 \IC(A)\equiv r(\nu)\pmod2,
 \qquad
 \IC(A)\ge r(\nu)\ge1.
\]
Consequently,
\begin{equation}
 q^2-k\le
 \binom{k}{2}(q-1)-\frac6m\binom{k}{4}-\frac1m.
\label{eq:nucleus-out}
\end{equation}
\end{proposition}

\begin{proof}
The nucleus is a required point outside $A\cup\cC$, so relative
completeness gives $r(\nu)\ge1$; Lemma~\ref{lem:max-index} gives the upper
bound.  A line is tangent to $\cC$ exactly when it passes through $\nu$.
Thus the number of tangent $A$-secants is $r(\nu)$.  Reducing the sum of the
intersection sizes $|\ell\cap\cC|$ modulo two gives
$\IC(A)\equiv r(\nu)\pmod2$.  The remaining claims follow from
Corollary~\ref{cor:conic-bound}.
\end{proof}

\begin{corollary}[Universal even-characteristic loss]
\label{cor:even-loss}
Every $\cC$-complete arc in even characteristic satisfies
\[
 \IC(A)\ge1.
\]
Consequently, inequality~\eqref{eq:conic-complete} can always be strengthened
by replacing its final term by at least $-1/m$.
\end{corollary}

\begin{proof}
If the nucleus belongs to $A$, Proposition~\ref{prop:nucleus-in} gives
$\IC(A)\ge k-1\ge1$.  Otherwise Proposition~\ref{prop:nucleus-out} gives
$\IC(A)\ge r(\nu)\ge1$.
\end{proof}

These two cases are useful structural reductions, but
Remark~\ref{rem:scale} limits what the incidence term alone can accomplish.
For example, at $(q,k)=(16,8)$ inequality~\eqref{eq:conic-complete} would
be contradicted only by $\IC(A)>268$, whereas the universal upper bound is
$\IC(A)\le56$.  Thus the exact determination of $\rhoC(16)$ below requires
information beyond a lower bound on $\IC(A)$ inserted into the present
inequality.

\section{Verified finite-field examples}
\label{sec:examples}

The verified values at $q=8,9,11$ use the lower bound $L_2(q)$ and explicit
coordinates.  At $q=16$, the same ingredients give $8\le\rhoC(16)\le9$;
a checked projective classification excludes eight.

The algebraic rejection used in that classification is an instance of the
following elementary linear-algebra observation.  Linear systems of
quadrics containing arcs are an established tool; see Glynn~\cite{Glynn1994}
and Ball~\cite{Ball2017Quadrics}.  Fix nonzero vector representatives of the
projective points under consideration.  If $V$ is a finite-dimensional
linear system of homogeneous forms, write
$\operatorname{ev}_x\in V^*$ for evaluation at the representative of $x$.
The zero or nonzero status of this evaluation is independent of the chosen
representative.

\begin{lemma}[Uncovered evaluation obstruction]
\label{lem:evaluation-obstruction}
Let $A$ and $U$ be finite sets of projective points.  There is no nonzero
$f\in V$ such that $U\subseteq Z(f)$ and $A\cap Z(f)=\varnothing$ if either
of the following holds:
\begin{enumerate}
 \item the evaluation map
 $f\mapsto(\operatorname{ev}_u(f))_{u\in U}$ is injective;
 \item for some $a\in A$, the functional $\operatorname{ev}_a$ belongs to
 the span of $\{\operatorname{ev}_u:u\in U\}$.
\end{enumerate}
\end{lemma}

\begin{proof}
If $f$ vanishes on $U$, injectivity in the first case gives $f=0$.  In the
second case, applying a spanning relation for $\operatorname{ev}_a$ to $f$
gives $f(a)=0$, contrary to $A\cap Z(f)=\varnothing$.
\end{proof}

For the full system of degree-$d$ forms, the evaluation functionals are the
degree-$d$ Veronese images of the points, up to nonzero scalar.  Thus the
first alternative says that the uncovered Veronese images span the dual
coefficient space; the second supplies a selected Veronese image already in
their span.  In particular, for the full degree-$d$ system it is enough to
exhibit $\binom{d+2}{2}$ uncovered points whose Veronese evaluation matrix is
invertible.  This formulation is not specific to conics.

The following sharp form is useful when a zero locus must avoid several
selected points simultaneously.  The vector-space covering threshold is
classical; see, for example, Khare~\cite{Khare2009}.

\begin{proposition}[Finite-field evaluation dichotomy]
\label{prop:evaluation-dichotomy}
Let $W$ be a finite-dimensional vector space over $\F_q$, let
$\nu:X\to W$ be any feature map, and let $U,A\subseteq X$ be finite with
$|A|\le q$.  There is a linear form $f\in W^*$ such that
\[
 f(\nu(u))=0\quad(u\in U),\qquad
 f\ne0,\qquad f(\nu(a))\ne0\quad(a\in A)
\]
if and only if
\[
 \langle\nu(U)\rangle\ne W
 \quad\hbox{and}\quad
 \nu(a)\notin\langle\nu(U)\rangle\quad(a\in A).
\]
The uniform threshold is sharp: $q+1$ distinct hyperplanes through a common
codimension-two subspace can cover $W$.
\end{proposition}

\begin{proof}
Put $S=\langle\nu(U)\rangle$.  Necessity follows by applying $f$ to the
stated span relations.  Conversely, $S\ne W$ makes the annihilator
$S^\perp\subseteq W^*$ nonzero.  For each $a\in A$, the condition
$\nu(a)\notin S$ says that evaluation at $\nu(a)$ cuts out a proper
hyperplane of $S^\perp$.  At most $q$ proper hyperplanes cannot cover a
nonzero vector space over $\F_q$: their union has size at most
$1+q(|S^\perp|/q-1)<|S^\perp|$.  Any form outside their union has the
required properties.  For sharpness, quotient by a common codimension-two
subspace and use all $q+1$ lines in $\F_q^2$.
\end{proof}

Taking $\nu$ to be the degree-$d$ Veronese map gives the same exact
dichotomy for homogeneous degree-$d$ forms.  Thus
Lemma~\ref{lem:evaluation-obstruction} is the one-selected-point exclusion
mechanism, while Proposition~\ref{prop:evaluation-dichotomy} determines
exactly when a single form can avoid as many as $q$ selected points.

\begin{proposition}\label{prop:small-values}
For the standard conic
\[
 \cC:\ XZ=Y^2
\]
in the indicated projective planes,
\[
 \rhoC(8)=6,
 \qquad
 \rhoC(9)=6,
 \qquad
 \rhoC(11)=6.
\]
\end{proposition}

\begin{proof}
Equation~\eqref{eq:L2} gives
\[
 L_2(8)=L_2(9)=L_2(11)=6.
\]
Appendix~\ref{app:witnesses} lists $\cC$-complete arcs of size six in each
plane.  Their verification is described below.
\end{proof}

\begin{theorem}[Quadratic avoidance over $\F_{16}$]
\label{thm:q16-quadratic}
Let $A$ be any eight-arc in $\PG(2,16)$, and let $U(A)$ be the points outside
$A$ that lie on no secant of $A$.  There is no nonzero homogeneous quadratic
form $Q$ such that
\[
 U(A)\subseteq Z(Q)
 \qquad\hbox{and}\qquad
 A\cap Z(Q)=\varnothing.
\]
The assertion includes singular quadratic forms.
\end{theorem}

\begin{proof}
Any four points of $A$ form a projective frame.  Normalize that frame and
apply canonical augmentation.  This gives respectively $4$, $61$, $454$,
and $2633$ projective classes at sizes five through eight (from $182$, $532$,
$5155$, and $21495$ legal extensions).  Every retained transition is
certified by an invertible $3\times3$ matrix together with pointwise
projective scalar equalities, so the classification does not trust the
canonical-labeling program.  The final count of $2633$ projective eight-arc
classes independently reproduces Theorem~3.8.1 of
Al-Seraji--Al-Ogali~\cite{AlSerajiAlOgali2018}; the count itself is not new.
More specifically, each \texttt{StepBook.coverage} theorem says that every
legal one-point extension of its parent is represented by a certified entry,
and \texttt{StepBooksValid} says that the books cover the current parent list
in exactly its listed order.  Every entry carries a checked projective
equivalence to a member of the next list.  The theorem
\texttt{classifiedAt\_level8\_of\_frame} composes those equalities, and the
frame-reduction theorem transports an arbitrary eight-arc into a listed
leaf.  Thus closure and exhaustiveness of the four class lists, not merely
the arithmetic of each listed member, are kernel-checked; the C++ generator
is not trusted for completeness.

For $2630$ of the classified leaves, six points of $U(A)$ give an invertible
$6\times6$ evaluation matrix, giving the first alternative of
Lemma~\ref{lem:evaluation-obstruction}.  In each of
the remaining three leaves the evaluation rows have rank five, but an
explicit linear combination places the evaluation functional of a point of
$A$ in their span, giving the second alternative.

Finally, the obstruction is projectively invariant.  If $gA$ is a normalized
representative, replace $Q$ by $Q\circ g^{-1}$.  This is nonzero because $g$
is invertible, its zero set is $gZ(Q)$, and projectivities carry secants to
secants, so $U(gA)=gU(A)$.  Thus a quadratic as in the statement would
contradict the checked obstruction for the normalized leaf.  The companion
Lean development kernel-checks every local alternative; its global conic
corollary additionally checks frame reduction and projective transport.
\end{proof}

\begin{corollary}[Exact relative value over $\F_{16}$]
\label{cor:q16-exact}
For every nonsingular conic $\cC\subseteq\PG(2,16)$,
\[
 \rhoC(16)=9.
\]
\end{corollary}

\begin{proof}
Equation~\eqref{eq:L2} gives $L_2(16)=8$.  If a $\cC$-complete eight-arc
existed, its ordinary-uncovered locus would lie in the zero set of the
nonzero quadratic defining $\cC$, while the arc is disjoint from that zero
set.  This contradicts Theorem~\ref{thm:q16-quadratic}.  The nine-point
witness in Appendix~\ref{app:witnesses} gives the matching upper bound.
\end{proof}

The published ordinary classification records arc counts, stabilizers,
secant-index distributions, and conic-contained classes.  We did not find
the uncovered-locus rank split, Theorem~\ref{thm:q16-quadratic}, or its exact
relative-conic corollary tabulated there.

\begin{remark}[Anatomy of the three exceptional leaves]
Using the binary $\F_{16}$ encoding from Appendix~\ref{app:witnesses} and
coefficient order $(X^2,Y^2,Z^2,XY,XZ,YZ)$, the one-dimensional kernels of
the three exceptional evaluation matrices are generated by
\[
 (1,1,1,1,1,0),\qquad (0,0,0,5,4,1),\qquad (2,1,1,5,5,1).
\]
The first and third forms split over $\F_{16}$ as
\[
 (X+6Y+6Z)(X+7Y+7Z)
 \quad\hbox{and}\quad
 2(X+4Y+15Z)(X+15Y+4Z),
\]
respectively, and each union of lines contains two selected points.  The
middle form is nonsingular and contains seven of the eight selected points.
Thus the exceptional leaves are not merely unexplained rank defects: two are
reducible quadratic configurations, while the third is an arc differing by
one point from a seven-point subset of a conic.  These descriptions follow
by direct field arithmetic and are not used to establish classification
completeness.
\end{remark}

For reference, the verified incidence statistics are
\[
\begin{array}{c|rrrrr}
q&|A|&\binom{|A|}{2}&q^2-|A|&\IC(A)&\max_{x\notin A}r(x)\\ \hline
8 &6&15&58 &16&3\\
9 &6&15&75 &13&3\\
11&6&15&115&0 &3\\
16&9&36&247&32&4
\end{array}
\]

The row at $q=11$ has a stronger geometric interpretation.  Since
$\IC(A)=0$, none of the 15 secants meets $\cC$: every secant of the witness
is exterior to the conic.  If $N_i$ denotes the number of required points
of index $i$, relative completeness and the maximum-index bound give only
$i=1,2,3$, while the two moment equations give
\[
 N_1+N_2+N_3=115,\qquad
 N_1+2N_2+3N_3=150,\qquad
 N_2+3N_3=45.
\]
Consequently $(N_1,N_2,N_3)=(90,15,10)$.

\begin{proposition}[The $q=11$ code and extension spectrum]
\label{prop:q11-code}
Let $A$ be the six-point witness in Appendix~\ref{app:witnesses}, and let
$C_A$ be its parity-check code.  Then:
\begin{enumerate}[label=(\roman*),leftmargin=2em]
 \item $C_A$ is a projectively non-GRS $[6,3,4]_{11}$ MDS code of covering
       radius three, and its projective deep-hole syndrome locus is exactly
       the twelve-point conic $\cC$;
 \item the affine cosets of distances $0,1,2,3$ number
       $(1,60,1150,120)$, and the distance-two cosets split as
       $(900,150,100)$ according as they have one, two, or three
       minimum-weight leaders;
 \item the simultaneous one-column MDS extensions form the independence
       complex of the icosahedral graph, with independence polynomial
       \[
        1+12t+36t^2+20t^3;
       \]
       exactly six two-column and twenty three-column extensions are
       maximal, giving six complete eight-arcs and twenty complete nine-arcs
       containing the fixed seed, and there is no ten-arc extension; and
 \item each witness column colours five pairwise disjoint edges of the
       icosahedron, missing an antipodal pair, and the six colour classes
       partition its thirty edges.  Adding the six antipodal edges turns
       them into six perfect matchings of the augmented graph.
\end{enumerate}
\end{proposition}

\begin{proof}
Every witness triple has nonzero determinant and the six columns span
$\F_{11}^3$, giving the MDS parameters.  The $6\times6$ evaluation matrix
of the quadratic monomials
\[
 X^2,\ XY,\ XZ,\ Y^2,\ YZ,\ Z^2
\]
is nonsingular, so no conic contains all six
projective columns; by the preceding definition and the classical
normal-rational-curve dictionary, the code is not projectively GRS.  The
formal theorem checks the no-conic premise; it does not build coding-theory
equivalences into the trusted kernel statement.  Direct finite
arithmetic gives the secant-index spectrum
\[
 0:12,\qquad 1:90,\qquad 2:15,\qquad 3:10,
\]
where index zero is precisely the conic.  Multiplying each nonzero
projective direction by its ten affine scalars gives all nonzero affine
syndromes bijectively.  In the companion formalization, membership in each
counted set is equivalent to actual parity-check distance, and for every
distance-two direction the determinant-zero witness pairs are proved equal
to the supports of actual weight-two coefficient words.  Uniqueness on an
independent support and scalar multiplication then give the claimed coset
and leader counts, not merely rescaled incidence counts.  The index-zero
directions give the covering radius.

For simultaneous extensions, two conic columns conflict exactly when their
chord contains a witness column.  No three conic columns are collinear, so
there is no additional triple obstruction and valid extension sets are
exactly graph-independent sets.  Exhausting the $2^{12}$ subsets gives the
displayed polynomial and maximal-set counts.  The same determinant table
partitions the thirty conflict edges into the six stated matchings.  All
these finite assertions are independently reproduced by the supplementary
Python and C++ checkers and kernel-checked in the companion Lean module.
Lean additionally proves that residual maximality upgrades to ordinary
completeness: relative completeness confines every further point to the
conic, and maximality excludes the remaining conic points.  It also checks
that the six missing antipodal edges are distinct and that adjoining the
appropriate edge to each colour class gives a six-edge perfect matching;
the six augmented classes are pairwise edge-disjoint and partition the
augmented graph.
\end{proof}

The six-set is a classical Clebsch hexagon with icosahedral $A_5$ symmetry;
this identification and the associated orbit geometry belong to prior art,
not to the proposition's novelty
posture~\cite{Edge1956,Dye1991,StormeVanMaldeghem1995}.  More precisely,
Dye's edge criterion at $q=11$ and the complete-exterior-set formulation of
Blokhuis--Seress--Wilbrink give $\cC(\F_{11})\subseteq U(A)$
\cite[discussion preceding Theorem~6, p.~281]{Dye1991}
\cite[pp.~143, 146]{BlokhuisSeressWilbrink1992}.  Dye's ten Brianchon
concurrences and the chord-defect identity then give $|U(A)|=12$, hence
equality.  This exact uncovered-locus conclusion is a short synthesis of
classical inputs rather than a new classical configuration theorem.
The exact conjunction of the relative-conic deep-hole locus, refined coset
counts, and extension polynomial was not located in the bounded search, so
we report it only as a checked synthesis rather than a priority claim.

\begin{remark}[Auxiliary residual-game structure]
There is also a game-theoretic consequence of this exceptional row, not used
in any bound above.  After the six witness points are occupied, all twelve
points of $\cC$ remain legal.  Join two conic points when their chord contains
a witness point.  The resulting $12$-vertex, $30$-edge, $5$-regular conflict
graph is the icosahedral graph.  Its antipodal involution is fixed-point-free,
preserves adjacency, and pairs only nonadjacent vertices, so the usual
copycat strategy makes the corresponding independent-set building position
a P-position.  The companion Lean development checks every seeded continuation,
proves an exact recursive localization from the projective cap game to this
independent-set game, and hence proves that the displayed six-point projective
position itself is P.  It also checks that the displayed $q=9$ six-arc covers
every projective point, including the conic, so that position is terminal P.
These observations are consequences of the displayed witnesses, not novelty claims.
\end{remark}

The supplementary Python program
\texttt{verify\_relative\_conic\_arcs.py} performs exact finite-field
arithmetic and, for each witness:

\begin{enumerate}[label=(\alph*),leftmargin=2em]
\item enumerates and normalizes all $q^2+q+1$ projective points;
\item checks that $XZ=Y^2$ contains exactly $q+1$ points;
\item verifies disjointness from the conic and that every pair-line contains
      exactly two selected points;
\item verifies that every point outside $A\cup\cC$ lies on a selected
      pair-line; and
\item independently checks both equations in
      Proposition~\ref{prop:moments}.
\end{enumerate}

Its SHA-256 digest is
\begin{center}
\texttt{e9508958d604e68c6c3d09fd3afadfaa8a3126508a51f1dfa993e7a7aed5d36a}.
\end{center}
The verifier is deliberately simpler than a search program: it certifies the
listed upper-bound witnesses but makes no claim about search completeness.

Two further implementations independently regenerate
Proposition~\ref{prop:q11-code}.  Their separately written source files are
\begin{center}
\texttt{check\_q11\_structure.py},\qquad
\texttt{check\_q11\_structure.cpp}.
\end{center}
Their SHA-256 digests are
\begin{center}\scriptsize
\texttt{0abe909c9aadce0db4c75f296c8de25e929dd1065c906996da8dec017e534d69},\\
\texttt{1753674172d48f1d056d350e30baa9eb67de0810c84a96da0440947768ae041c}.
\end{center}
They agree on the projective, syndrome, chord, and group-action data.  Both
include transformed-coordinate invariance checks; the adversarial controls
replace one witness point and mutate a group generator, respectively, and
reject the exceptional conclusions.  The small-field hyperplane checker
\texttt{check\_evaluation\_dichotomy.py}, with digest
\begin{center}
\texttt{6ed309bd2461ce9998cbd3bcaa5396379e6973b7503ce2dcdfb32c9806386566},
\end{center}
exhausts every family of at most $q$ projective hyperplanes for
$q=2,3,5$ and dimensions one through three, and checks the sharp $q+1$
plane cover.

The separate exact-classification generator
\texttt{search\_rhoc16.cpp} has SHA-256 digest
\begin{center}
\texttt{589af8430e94b4c9c23ce895e6d32d2b3b5b9b387b1fb23ed6d3875cdee39031}.
\end{center}
Its frozen report \texttt{search\_rhoc16\_output.txt} has digest
\begin{center}
\texttt{6989079b5cb64b0e57d5c42b872093fff99f861300b8fbb909daef450c15cc63}.
\end{center}
The generator emits local Lean certificates; its own assertions and
canonical labels are not part of the trusted proof boundary.

The companion Lean library \texttt{lean/RelativeConicArcs/} formalizes the
incidence, defect, conic, asymptotic, averaging, nucleus, evaluation-dichotomy,
syndrome, and coding arguments.  The downstream \texttt{Q11Coding} module
kernel-checks the exact code parameters, covering radius, deep-hole locus,
leader distributions, quadratic obstruction, chord partition, and extension
complex.  Its
generic Boolean checker inspects conic disjointness, the arc condition, and
coverage of the canonical projective representatives
$[1:y:z]$, $[0:1:z]$, and $[0:0:1]$; a proved normalization and soundness
theorem transports acceptance to relative completeness.  The concrete
checks use kernel reduction rather than \texttt{native\_decide}.  The trust
manifest \texttt{lean/RelativeConicArcs/TRUST.md} records the verifier hash,
the theorem map, and axiom profiles.  In particular, the finite results do
not depend on the external Kim--Vu input used for the conditional asymptotic
upper bound.

\section{Further questions}

The elementary theory leaves several concrete problems.

\begin{problem}
Determine whether the uncovered-quadratic-rank obstruction used at $q=16$
extends to other even orders or gives useful lower bounds on the size of a
$\cC$-complete arc beyond the scalar defect inequality.
\end{problem}

\begin{problem}
Develop the higher-dimensional analogue for caps complete outside a
prescribed quadric in $\PG(n,q)$.  The evaluation-avoidance lemma already
applies to arbitrary feature maps and every Veronese degree; the missing
ingredients are the appropriate cap secant moments and useful examples in
dimensions at least three.
\end{problem}

\begin{problem}
Determine whether the exact defect remainder yields new lower bounds for
arc-constrained saturating sets or the corresponding covering codes, and
whether small-order scans beyond $q=11$ produce further recognizable
extension complexes.  The evaluation formulation is also the elementary
linear-algebraic interface used by polynomial-method arguments, but no
slice-rank or cap-set consequence is asserted here.
\end{problem}

\begin{problem}
Construct $\cC$-complete arcs of size $O(\sqrt q)$, or prove that this is
impossible for an infinite family.  Theorem~\ref{thm:transfer} currently
supplies only the general complete-arc bound
$O(\sqrt q(\log q)^c)$.
\end{problem}

\begin{problem}
Classify equality and bounded-defect configurations in
Theorem~\ref{thm:defect}.  The equality pattern requires all nontrivial
secant concurrence to occur at points of maximum possible index, a severe
combinatorial restriction.
\end{problem}

\begin{problem}
Determine whether the conic stabilizer or its coherent configuration yields
useful finite-order restrictions not reducible to a bound on $\IC(A)$.
Remark~\ref{rem:scale} shows that an estimate for $\IC(A)$ alone cannot
improve the asymptotic constant obtained here.
\end{problem}

The central distinction is between universal overlap, already measured by
the classical second index equation, and genuinely conic-specific geometry.
The former gives the lower bound in this paper.  Progress beyond it will
probably require an explicit construction, an algebraic obstruction, or a
classification of near-extremal secant arrangements rather than another
application of the same incidence budget.

\appendix
\section{Witness coordinates and field encodings}
\label{app:witnesses}

All coordinates are homogeneous and use the standard conic $XZ=Y^2$.

For $q=8$, write
$\F_8=\F_2[\alpha]/(\alpha^3+\alpha+1)$ and encode
$a_0+a_1\alpha+a_2\alpha^2$ by the binary integer
$a_0+2a_1+4a_2$.  A witness is
\[
\begin{split}
 &(0,1,1),(0,1,2),(1,0,1),\\
 &(1,0,2),(1,1,2),(1,1,6).
\end{split}
\]

For $q=9$, write
$\F_9=\F_3[\alpha]/(\alpha^2+1)$ and encode
$a_0+a_1\alpha$ by $a_0+3a_1$.  A witness is
\[
\begin{split}
 &(1,0,4),(1,0,5),(1,1,0),\\
 &(1,1,2),(1,2,3),(1,2,4).
\end{split}
\]

For $q=11$, coordinates lie in the prime field.  A witness is
\[
\begin{split}
 &(1,10,0),(1,9,1),(1,4,7),\\
 &(1,8,5),(0,1,4),(1,1,7).
\end{split}
\]

For $q=16$, write
$\F_{16}=\F_2[\alpha]/(\alpha^4+\alpha+1)$ and use binary coefficient
encoding.  A witness of size nine is
\[
\begin{split}
 &(1,5,14),(1,5,3),(1,9,2),(1,0,9),(1,13,13),\\
 &(1,6,10),(1,3,9),(0,1,11),(1,10,2).
\end{split}
\]

\begin{thebibliography}{99}

\bibitem{AlabdullahHirschfeld2019}
S.~Alabdullah and J.~W.~P.~Hirschfeld,
\newblock A new lower bound for the smallest complete $(k,n)$-arc in
$\PG(2,q)$,
\newblock \emph{Designs, Codes and Cryptography} \textbf{87} (2019), 679--683.
\newblock doi:10.1007/s10623-018-00592-8.

\bibitem{AlSerajiAlOgali2018}
N.~A.~M. Al-Seraji and R.~A.~B. Al-Ogali,
\newblock Classification of arcs in finite projective plane of order sixteen,
\newblock \emph{Al-Mustansiriyah Journal of Science} \textbf{29}(1) (2018),
138--149.
\newblock doi:10.23851/mjs.v29i1.184.

\bibitem{Ball1997}
S.~Ball,
\newblock On small complete arcs in a finite plane,
\newblock \emph{Discrete Mathematics} \textbf{174} (1997), 29--34.
\newblock doi:10.1016/S0012-365X(96)00315-9.

\bibitem{Ball2017Quadrics}
S.~Ball,
\newblock On arcs and quadrics,
\newblock in \emph{Arithmetic of Finite Fields---WAIFI 2016}, Lecture Notes
in Computer Science \textbf{10064}, Springer, 2017, 95--102.
\newblock doi:10.1007/978-3-319-55227-9\_7.

\bibitem{BartoliEtAl2018}
D.~Bartoli, A.~A.~Davydov, S.~Marcugini, and F.~Pambianco,
\newblock On the smallest size of an almost complete subset of a conic in
$\PG(2,q)$ and extendability of Reed--Solomon codes,
\newblock \emph{Problems of Information Transmission} \textbf{54} (2018),
101--115.
\newblock doi:10.1134/S0032946018020011.

\bibitem{BlokhuisSeressWilbrink1992}
A.~Blokhuis, A.~Seress, and H.~A. Wilbrink,
\newblock Characterization of complete exterior sets of conics,
\newblock \emph{Combinatorica} \textbf{12}(2) (1992), 143--147.
\newblock doi:10.1007/BF01204717.

\bibitem{Edge1956}
W.~L.~Edge,
\newblock Conics and orthogonal projectivities in a finite plane,
\newblock \emph{Canadian Journal of Mathematics} \textbf{8} (1956), 362--382.

\bibitem{DavydovEtAl2021}
A.~A.~Davydov, S.~Marcugini, and F.~Pambianco,
\newblock On the weight distribution of the cosets of MDS codes,
\newblock \emph{Advances in Mathematics of Communications} \textbf{17}(5)
(2023), 1115--1138.
\newblock doi:10.3934/amc.2021042.
\newblock arXiv:2101.12722v2.
\newblock (Numbered results are cited by their arXiv~v2 numbering.)

\bibitem{Dye1991}
R.~H. Dye,
\newblock Hexagons, conics, $A_5$ and $\operatorname{PSL}_2(K)$,
\newblock \emph{Journal of the London Mathematical Society} (2)
\textbf{44} (1991), 270--286.
\newblock doi:10.1112/jlms/s2-44.2.270.

\bibitem{Glynn1994}
D.~G. Glynn,
\newblock On the construction of arcs using quadrics,
\newblock \emph{Australasian Journal of Combinatorics} \textbf{9} (1994),
3--19.

\bibitem{Hirschfeld1998}
J.~W.~P.~Hirschfeld,
\newblock \emph{Projective Geometries over Finite Fields}, second edition,
\newblock Oxford University Press, 1998.

\bibitem{KimVu2003}
J.~H.~Kim and V.~H.~Vu,
\newblock Small complete arcs in projective planes,
\newblock \emph{Combinatorica} \textbf{23} (2003), 311--363.
\newblock doi:10.1007/s00493-003-0024-1.

\bibitem{Kaipa2017}
K.~Kaipa,
\newblock Deep holes and MDS extensions of Reed--Solomon codes,
\newblock \emph{IEEE Transactions on Information Theory} \textbf{63}
(2017), 4940--4948.

\bibitem{Khare2009}
A.~Khare,
\newblock Vector spaces as unions of proper subspaces,
\newblock \emph{Linear Algebra and its Applications} \textbf{431} (2009),
1681--1686.
\newblock doi:10.1016/j.laa.2009.06.001.

\bibitem{Nagy2018}
Z.~L.~Nagy,
\newblock Saturating sets in projective planes and hypergraph covers,
\newblock \emph{Discrete Mathematics} \textbf{341} (2018), 1078--1083.
\newblock doi:10.1016/j.disc.2018.01.011.

\bibitem{NgWild2001}
S.-L. Ng and P.~R. Wild,
\newblock On $k$-arcs covering a line in finite projective planes,
\newblock \emph{Ars Combinatoria} \textbf{58} (2001), 289--300.

\bibitem{StormeVanMaldeghem1995}
L.~Storme and H.~Van Maldeghem,
\newblock Primitive arcs in $\PG(2,q)$,
\newblock \emph{Journal of Combinatorial Theory, Series A} \textbf{69}
(1995), 200--216.
\newblock doi:10.1016/0097-3165(95)90051-9.

\bibitem{StormeThas1991}
L.~Storme and J.~A. Thas,
\newblock Generalized Reed--Solomon codes and normal rational curves: an
improvement of results by Seroussi and Roth,
\newblock in \emph{Advances in Finite Geometries and Designs}, Oxford
University Press, 1991, 369--389.
\newblock doi:10.1093/oso/9780198535928.003.0031.

\end{thebibliography}

\end{document}
